% ============================================================
%  ABSTRACT
% ============================================================

\cleardoublepage
\thispagestyle{plain}

\begin{center}
  {\Large\scshape Abstract}
\end{center}

\vspace{1.5em}

\noindent
This monograph develops a categorical theory of thermodynamics from a single axiom
about bounded phase space. Three structural consequences of the axiom~---
boundedness, the No-Null condition that no persistent state has zero dynamical
activity, and finite observational resolution~--- generate a framework in which
entropy is the logarithm of the categorical count of distinguishable states accessible
to a finite observer through partition operations.

\medskip

The triple equivalence theorem establishes that oscillatory dynamics, categorical
state enumeration, and partition arithmetic yield identical entropy
$S = k_{\mathrm{B}} M \ln n$, with $W = n^M$ in each formulation.
Boltzmann's formula stands at the foundation; the categorical extension applies it
past regimes where the kinetic interpretation runs out. The framework dissolves
three paradoxes that have stood as the field's most enduring open problems:
Gibbs's paradox of mixing, resolved through observer-relative distinguishability;
Maxwell's Demon, resolved through the heat-entropy decoupling that separates
statistically emergent quantities from categorical fundamentals; and Kelvin's heat
death, resolved through the kinetic-to-categorical phase transition.

\medskip

Counting from the heat-death configuration of the observable universe~---
approximately $10^{80}$ particles distributed at maximum spatial separation, with
$10^{84}$ distinguishable single-particle states~--- the monograph derives the
maximum categorical count as the tetration
$N_{\max} \approx (10^{84})\uparrow\uparrow(10^{80})$, exceeding every previously
named large number by tetrational orders of magnitude. The kinetic phase of cosmic
evolution, in which conventional thermodynamics operates, is shown to occupy a finite
slice of this count; the categorical phase, beginning at heat death, exceeds the
kinetic phase by $N_{\max}$ orders of magnitude.

\medskip

Five thermodynamic regimes~--- ideal gas, plasma, degenerate matter, relativistic gas,
and Bose-Einstein condensate~--- emerge as limiting cases of one underlying partition
geometry. Classical thermodynamics is exhibited as the specialisation to sparse
partition networks in thermal equilibrium. The Poincar\'e duality between computation
and thermodynamics establishes that categorical state traversal and physical evolution
are two descriptions of the same partition operation. The mass spectrometer is
constructed as a categorical measurement instrument, with four hardware oscillators
forming a complete set of commuting observables on partition coordinate space.
Spectral-native dynamics for both gas and liquid regimes are developed on a common
GPU shader architecture, instantiating the rendering-measurement identity at the
implementation level.

\medskip

The final chapters establish that the universe is cyclic. After $N_{\max}$ categorical
steps, the enumeration completes; the only category that remains unfilled is the
singular state, structurally identical to the cosmological singularity from which
the cycle began. The cyclic structure is not posited as a cosmological proposal
but derived as a consequence of categorical completion. The straight line of cosmic
evolution, which the kinetic formalism reads as a one-way arrow from low entropy to
heat death, is exhibited as a closed curve whose local direction is forward at every
point and whose global structure closes.

\medskip

The monograph contains fifteen chapters across seven parts: the observation boundary,
the dissolution of the classical paradoxes, the categorical thermodynamic framework,
equations of state across all five regimes, classical thermodynamics as a special
case, the measurement instrument, and dynamics. Validation proceeds through
derivation rather than fitting; there are no tunable parameters and no adjustments
to data. The framework operates on fundamental constants only: $k_{\mathrm{B}}$,
$\hbar$, $c$, and the structural integers that follow from the partition geometry.

\vspace{2em}

\noindent
\textbf{Keywords:} bounded phase space, categorical entropy, triple equivalence,
partition operations, Second Law, Boltzmann formula, Gibbs paradox, Maxwell's Demon,
Kelvin heat death, Loschmidt paradox, equations of state, mass spectrometry,
spectral-native dynamics.

\cleardoublepage
